% Integration fragment. Author: GPT-6 Astra Pro, inherited effort.
% 21 November 2026. Source proof: CM_PERFECT_COMPLEX.md.
% The final HD-N06 coherence argument has its own check events.
\section{An explicit equivariantly semiregular complex at a CM sixfold}

We construct a perfect complex on the standard balanced principally
polarized sixfold
\[
 A_d=E_i^3\times E_i^3,\qquad
 \iota(i)=\operatorname{diag}(i_0,-i_0),\qquad
 \theta_0=H_++H_-,
\]
where $E_i=\mathbb C/(\mathbb Z+i\mathbb Z)$ and $i_0$ is its multiplication
by $i$. The construction supplies a coherent finite translation/tensor
action of order $4096$, negative self-Ext vanishing, and injectivity of
the full semiregularity trace on the invariant obstruction space. Its
entire character is
\begin{equation}\label{hd-n-final-character}
 \operatorname{ch}(\mathcal E_d)
 =32t_d+4-8\theta_0^2+\frac83\theta_0^4-\frac{16}{45}\theta_0^6,
 \qquad B_0=0,
\end{equation}
where $t_d$ is primitive integral, $\theta_0t_d=0$, and
$(-1)^3\int t_d^2=32$. Thus the class normalization is
$(\delta,T)=(1,32)$ and the Weil coefficient is $\lambda=32$.
Markman's supplied manuscript motivates the ideal-curve and Fourier--Mukai
design; no semiregularity or algebraicity theorem of that manuscript is
used as an established input.

\subsection{HD-N01: a Fermat quotient and a disjoint two-curve orbit}
\label{hd-n01}

Let $C\subset\mathbb P^2$ be the plane quartic
$x^4+y^4+z^4=0$. It is smooth, nonhyperelliptic, and has genus $3$.
The three sign involutions form a Klein four group. Each involution has
four fixed points, so its quotient has genus $1$ by Riemann--Hurwitz.
The quotient by the full Klein group is the conic
$X^2+Y^2+Z^2=0$ under $[x:y:z]\mapsto[x^2:y^2:z^2]$. Hence
$H^1(C,\mathbb Q)$ is the direct sum of its three nontrivial Klein
character spaces, each of dimension $2$.

For example, one elliptic quotient has equation $u^2=-(t^4+1)$.
The automorphism $(t,u)\mapsto(it,u)$ fixes a point at $t=0$ and acts by
$i$ on $dt/u$. With that fixed point as origin its Jacobian is $E_i$.
Write $f_j:C\to E_i$ for the three quotient maps and put
\[
 q:E_i^3\longrightarrow J(C),\qquad
 q(z_1,z_2,z_3)=\sum_{j=1}^3f_j^*z_j.
\]
The cross norms $f_{j,*}f_k^*$ vanish for $j\ne k$: they intertwine
different Klein characters, hence vanish rationally, and a homomorphism
of abelian varieties which vanishes rationally is zero. The diagonal
norms are $[2]$. If $H$ is the product principal polarization, this gives
\begin{equation}\label{hd-n-fermat-isogeny}
 q^*\Theta=2H,\qquad \deg q=8,\qquad \ker q\subset E_i^3[2].
\end{equation}
Indeed the polarization homomorphism of $q^*\Theta$ is diagonal $[2]$;
positivity makes $q$ an isogeny and
$\chi(q^*\Theta)=2^3=\deg q$. The norm identities also give
$[2]=q^\dagger q$, proving the kernel assertion.

Choose $0\ne\tau\in E_i[2]$ and set
$\epsilon=f_1^*\tau$. The homomorphism $f_1^*$ is injective. A nonzero
point in its kernel would be $2$-torsion; trivializing its pullback would
factor the degree-two map $C\to E_i$ through the corresponding connected
\'etale double cover of $E_i$. The remaining map would have degree one,
forcing $C$ to have genus one. Thus $\epsilon\ne0$ and $2\epsilon=0$.

For an Abel--Jacobi image $C_0\subset J(C)$, the curves $C_0$ and
$C_0+\epsilon$ are disjoint. An intersection would give
$p-q=\epsilon$, with $p\ne q$, and hence $2p\sim2q$; this supplies a
degree-two function and contradicts nonhyperellipticity. Therefore the
quotient
\[
 u:J(C)\longrightarrow Y=J(C)/\langle\epsilon\rangle
\]
embeds $C_0$ as a smooth curve $\overline C$. Define
\[
 r=uq:X=E_i^3\longrightarrow Y,\qquad H_0=\ker r.
\]
Since $\epsilon=q(\tau,0,0)$,
\begin{equation}\label{hd-n-H0}
 H_0=\ker q+\langle(\tau,0,0)\rangle\subset X[2],
 \qquad |H_0|=16.
\end{equation}

\subsection{HD-N02: the full trace for the quotient ideal}
\label{hd-n02}

Put $\overline F=I_{\overline C}$ on the abelian threefold $Y$, and let
$h\in\operatorname{NS}(Y)_\mathbb Q$ be characterized by $u^*h=\Theta$.
The Abel curve has class $\Theta^2/2$ and Euler characteristic $-2$.
Since $\operatorname{td}(J(C))=1$, Riemann--Roch gives
$\operatorname{ch}(\mathcal O_{C_0})=\Theta^2/2-2[\mathrm{pt}]$.
Pullback of $\overline F$ is the ideal of the disjoint pair of curves, so
\begin{equation}\label{hd-n-ideal-character}
 \operatorname{ch}(\overline F)=1-h^2+\frac23h^3.
\end{equation}

The restriction $H^1(\mathcal O_Y)\to H^1(\mathcal O_{\overline C})$
is an isomorphism: its pullback through $u$ is the Abel--Jacobi
isomorphism. The normal sequence gives
\[
 H^0(N_{\overline C/Y})/H^0(T_Y)
 =\ker\bigl(H^1(T_C)\longrightarrow H^1(T_Y|_C)\bigr).
\]
The last map is injective. Its dual contains the canonical
multiplication
\[
 \operatorname{Sym}^2H^0(K_C)\longrightarrow H^0(K_C^2),
\]
which is an isomorphism: both spaces have dimension six, and the plane
quartic has no quadratic equation. Thus $h^0(N_{\overline C/Y})=3$.

Apply $\operatorname{Hom}(\overline F,-)$ to the ideal sequence. The
map $\operatorname{Hom}(\overline F,\overline F)\to
\operatorname{Hom}(\overline F,\mathcal O_Y)$ is the isomorphism
$\mathbb C\to\mathbb C$, yielding
\[
 0\longrightarrow H^0(N_{\overline C/Y})
 \longrightarrow\operatorname{Ext}^1(\overline F,\overline F)
 \longrightarrow\operatorname{Ext}^1(\overline F,\mathcal O_Y).
\]
The cohomology sequence gives
$H^2(\overline F)\simeq H^2(\mathcal O_Y)$ of dimension three.
Serre duality therefore gives
$\dim\operatorname{Ext}^1(\overline F,\mathcal O_Y)=3$.
Translations give three independent determinant-preserving directions;
tensoring by $\operatorname{Pic}^0(Y)$ gives three independent
determinant directions. Consequently
\begin{equation}\label{hd-n-ext-dimensions}
 \operatorname{ext}^1(\overline F,\overline F)
 =\operatorname{ext}^2(\overline F,\overline F)=6.
\end{equation}

Let $\sigma^j_F:\operatorname{Ext}^j(F,F)\to HH_{-j}(Y)$ be the
categorical trace. Buchweitz--Flenner, \emph{Adv. Math.} 217 (2008),
243--281, Proposition 6.2.1, Corollary 6.3.2 and Theorem 6.4.2, identify
its composition with Hochschild evaluation with contraction against
$\operatorname{ch}(F)$ under HKR. The tangent bundle is trivial, so the
Todd correction is one.

In degree two the entire target is
\[
 HH_{-2}(Y)=H^2(\mathcal O_Y)\oplus H^3(\Omega_Y^1),
 \qquad\dim HH_{-2}(Y)=6.
\]
The $H^2(\mathcal O_Y)$ part of Hochschild cohomology maps identically
to the first block through the rank-one term. The $H^1(T_Y)$ part maps
onto the second block by $-\kappa\mathbin{\lrcorner}h^2$.
Indeed contraction with $h$ surjects onto $H^2(\mathcal O_Y)$ and
wedging with $h$ is an isomorphism from that space to
$H^3(\Omega_Y^1)$. In coordinates $h=\sum x_jy_j$, it sends
$y_1y_2$ to a nonzero multiple of $x_3y_1y_2y_3$, and similarly for the
other basis elements. The contraction/evaluation composition has rank
six; \eqref{hd-n-ext-dimensions} therefore proves that $\sigma^2$
is injective.

In degree one, evaluation
$HH^1(Y)=H^1(\mathcal O_Y)\oplus H^0(T_Y)\to
\operatorname{Ext}^1(\overline F,\overline F)$ is precisely the six
Picard and translation directions, hence an isomorphism. Its trace has
rank six: the Picard directions inject through the rank term into
$H^1(\mathcal O_Y)$, and the translation directions inject through
$-v\mathbin{\lrcorner}h^2$ into $H^2(\Omega_Y^1)$.
The latter injectivity follows from nondegeneracy of $h$ and injectivity
of wedging $h$ on $H^1(\mathcal O_Y)$. Thus $\sigma^1$ is injective;
also $\sigma^0(\mathrm{id})=\operatorname{ch}(\overline F)\ne0$.

\subsection{HD-N03: external product and a coherent finite action}
\label{hd-n03}

The product $\overline F\boxtimes\overline F$ is semiregular on
$Y\times Y$. K\"unneth decomposes its second self-Ext into
\[
 \operatorname{Ext}^2\otimes\mathbb C,
 \qquad\operatorname{Ext}^1\otimes\operatorname{Ext}^1,
 \qquad\mathbb C\otimes\operatorname{Ext}^2.
\]
The trace maps these to the distinct summands
$HH_{-2}\otimes HH_0$, $HH_{-1}\otimes HH_{-1}$ and
$HH_0\otimes HH_{-2}$, by
$\sigma^2\otimes\operatorname{ch}$,
$\sigma^1\otimes\sigma^1$, and
$\operatorname{ch}\otimes\sigma^2$, with the graded tensor sign.
This follows directly from additivity of the Atiyah class for external
products and multiplicativity of trace. All factors are injective by
HD-N02.

Let $F_0=r^*\overline F$. Its canonical descent linearization gives
\[
 \operatorname{Ext}^j(F_0,F_0)^{H_0}
 \simeq\operatorname{Ext}^j(\overline F,\overline F).
\]
The pullback on Hochschild homology is an isomorphism, since $r$ is an
isogeny. Naturality of trace makes $F_0\boxtimes F_0$ strongly
$H_0^2$-semiregular, with the actual coherent descent action.

Let $L=\mathcal O_X(2H)$ and $\mu(x,y)=(x+y,y)$. Write $\mathcal P$
for the normalized Poincar\'e bundle and define
\begin{equation}\label{hd-n-perfect-complex}
 \Phi=(\mathrm{id}_X\boxtimes\mathrm{FM}_{\mathcal P^{-1}[3]})
        \circ\mu^*,\qquad
 \mathcal E=\Phi\bigl((F_0\otimes L)\boxtimes(F_0\otimes L)\bigr).
\end{equation}
This is a perfect complex. Mukai, \emph{Nagoya Math. J.} 81 (1981),
153--175, Theorem 2.2, p.~156, proves the Fourier--Mukai equivalence and
the inverse used here. Formula (3.1), p.~158, exchanges translations and
tensoring by degree-zero line bundles. The other factors are an
automorphism pullback and a line-bundle tensor. In particular
$\operatorname{Ext}^{<0}(\mathcal E,\mathcal E)=0$.

Transport the complete coherent $H_0^2$ action through these equivalences.
For $(x,y)\in H_0^2$ the resulting pair is translation by $(x+y,0)$
and tensoring by $p_2^*P_y$: the possible signs disappear since every
point involved is $2$-torsion and $\phi_L=[2]$. The map
$(x,y)\mapsto(x+y,y)$ is injective. Hence the transported action is by
an actual finite subgroup
\begin{equation}\label{hd-n-G}
 G=\{((a,0),(0,b)):a,b\in H_0\}\subset
 (X\times\widehat X)\times(\widehat X\times X),\qquad |G|=256.
\end{equation}
The coherence is inherited from the original descent action, not
inferred merely from this notation for its points. Functoriality of
trace under equivalences proves injectivity on
$\operatorname{Ext}^2(\mathcal E,\mathcal E)^G$. Thus $\mathcal E$ is
strongly $G$-semiregular, and weakly so by Perry, Definition 2.6 and
Lemma 2.7.

\subsection{HD-N04: the exact whole Chern character}
\label{hd-n04}

Choose integral classes $a_j,b_j,u_j,v_j$ on the three elliptic pairs,
with $\int a_jb_j=\int u_jv_j=1$. Put
\[
 H_x=\sum a_jb_j,\quad H_{\widehat x}=\sum u_jv_j,\quad
 P_j=a_jv_j-b_ju_j,\quad P=\sum P_j,\quad
 \theta=4H_x+H_{\widehat x},\quad B=-P/2.
\]
The integral polarization $\theta$ has type $(1,1,1,4,4,4)$ and volume
$64$. For $s\in\{i,-i\}$ define
\[
 \Omega_{s,j}=4a_jb_j+2sP_j-u_jv_j,\qquad
 W_s=\prod_{j=1}^3\Omega_{s,j},\qquad
 t=-\frac12(W_i+W_{-i}).
\]
We claim the exact full identity
\begin{equation}\label{hd-n-original-character}
 e^B\operatorname{ch}(\mathcal E)
 =t+64-2\theta^2+\frac{\theta^4}{96}-\frac{\theta^6}{46080}.
\end{equation}
To calculate it, note $r^*h=2H$ and expand
\[
 \operatorname{ch}(F_0\otimes L)
 =1+2H-2H^2-\frac43H^3
 =\alpha e^{2iH}+\beta e^{-2iH},\qquad
 \alpha=\frac{1-i}{2},\quad\beta=\frac{1+i}{2}.
\]
Grothendieck--Riemann--Roch for $\Phi$ uses the kernel character
$e^{-P}$; the shift $[3]$ contributes a minus sign and every Todd
class is one. On one elliptic pair, with source variables $a,b,c,d$
and output variables $a,b,u,v$, the required fibre integral is
\[
 \int_{c,d}\exp\bigl(2s(a+c)(b+d)+2t'cd-(cv-du)\bigr).
\]
For $s=t'\in\{i,-i\}$ it is
\[
 4s\exp\bigl(P_{\rm out}/2+s(2ab+uv/2)/2\bigr).
\]
For $t'=-s$, multiplication by $e^{-P_{\rm out}/2}$ gives
$4ab+2sP_{\rm out}-uv=\Omega_s$.
These four identities follow by finite expansion in the six-generator
exterior algebra. For the opposite signs one can also use
$\Omega_s=-(u-2sa)(v-2sb)$, annihilated by $2sab-P_{\rm out}/2$.
Multiplying the three coordinate factors, including the shift sign and
$\alpha^2,\beta^2,\alpha\beta$, yields
\[
 32e^{i\theta/4}+32e^{-i\theta/4}-\frac12(W_i+W_{-i}),
\]
whose expansion is \eqref{hd-n-original-character}.
The file \texttt{cm\_complex\_character.py} independently expands
these exterior algebras over $\mathbb Q(i)$ and checks the identities
exactly. Odd occurrences of $i$ cancel in $W_i+W_{-i}$, and all even
terms have even integral coefficients, so $t$ is integral.

\subsection{HD-N05: the Weil action and the exact balanced isogeny}
\label{hd-n05}

On each elliptic pair define the rational action of $i$ by
\[
 J=\begin{pmatrix}0&1/2\\-2&0\end{pmatrix}.
\]
It squares to $-1$, and has holomorphic eigenspace dimensions $(3,3)$.
For the rational polarization $\Theta=2H_x+H_{\widehat x}/2$ its
Rosati adjoint is $-J$; $\theta=2\Theta$ is the integral polarization
above, and the order $\mathbb Z[2i]$ acts by actual endomorphisms.
The six forms $u_j-2sa_j,v_j-2sb_j$ span one $J$-eigenspace, and
\[
 \Omega_{s,j}=-(u_j-2sa_j)(v_j-2sb_j)
\]
places $W_s$ on its Weil determinant line. Both conjugate lines are
nonzero, so $t\ne0$ and is a rational Weil class of type $(3,3)$.
For each pair,
\[
 \Omega_{i,j}\Omega_{-i,j}=-16a_jb_ju_jv_j,\qquad
 (4a_jb_j+u_jv_j)\Omega_{s,j}=0.
\]
It follows that $\theta t=0$ and $\int t^2=-2048$, so the displayed
action has normalization $(\delta,T)=(1,2048)$.

The actual isogeny from the standard diagonal balanced $A_d$ is, on
each pair,
\begin{equation}\label{hd-n-S}
 S=\begin{pmatrix}1&1\\2i_0&-2i_0\end{pmatrix},\qquad
 JS=S\operatorname{diag}(i_0,-i_0),\qquad
 \deg S=16^3=4096.
\end{equation}
Direct matrix multiplication gives
\[
 S^*(2H_x+H_{\widehat x}/2)=4(H_++H_-).
\]
These are the precise morphism and polarization comparison between the
two actions. The finite-action argument required after this isogeny is
the next finding.

\subsection{HD-N06: a coherent character summand on the standard CM point}
\label{hd-n06}

Relabel $G=H_0\times H_0$ by its translation and tensor coordinates
$(a,b)$, and let $K=\{0\}\times H_0$ be its pure-tensor subgroup.
We first identify its categorical coherence.

Put $Q=\mathcal O_X(H)$, with its symmetric product theta structure,
so $L=Q^2$. The theorem of the square gives the exact normalized
line-bundle identity
\begin{align}
 M:=\mu^*(L\boxtimes L)
 &=p_1^*Q^2\otimes p_2^*Q^4\otimes\mathcal P_H^2\notag\\
 &=(\mathrm{id}_X\times[2])^*
   (p_1^*Q^2\otimes p_2^*Q\otimes\mathcal P_H).
 \label{hd-n-M-descent}
\end{align}
Thus $M$ has the canonical pullback linearization under all translations
$(0,b)$, $b\in X[2]$. The subgroup $K$ corresponds to diagonal
translations $(b,b)$ before $\mu$, and to $(0,b)$ after it.
Equation \eqref{hd-n-M-descent} makes conjugation by tensor $M$
equivariantly identical to the canonical translation action on this
subgroup. The normalized Fourier--Mukai kernel carries it to the
canonical tensor action by $P_{-b}$. Hence this $K$-action has ordinary
coherent tensor multiplication, with no projective two-cocycle.

Let $\mathcal E'=S^*\mathcal E$, $C_S=\ker S$, and
\[
 T_S=S^{-1}(H_0\times\{0\})\subset A_d.
\]
Then $|C_S|=4096$ and $|T_S|=65536$. Pullback gives a coherent lifted
action of $T_S\times K$ on $\mathcal E'$ whose functor at $(z,b)$ is
translation by $z$, followed by tensoring with $S^*p_2^*P_b$.
Its map to the old group is $(z,b)\mapsto(S(z),b)$; its kernel $C_S$
acts by canonical finite-\'etale descent.

Each $S^*p_2^*P_b$ is trivial, since $b\in X[2]$ and the second row
of $S$ is $(2i_0,-2i_0)$. Compatible trivializations exist: the direct
sum with its tensor multiplication is a finite \'etale algebra with
all line summands trivial. On connected proper $A_d$ this is a constant
$K^\vee$-torsor; choose a point over the origin. Thus $K$ acts on
$\mathcal E'$ by commuting involutions. Its actual character
idempotents are
\[
 e_\chi=\frac1{16}\sum_{b\in K}\chi(b)^{-1}\rho(b).
\]
Idempotent completeness (Stacks, Lemma 21.47.8, Tag 08GA) gives
$\mathcal E'=\bigoplus_\chi F_\chi$, $F_\chi=e_\chi\mathcal E'$.

The deck group $C_S$ acts on the character labels through the descent
pairing
\[
 C_S\times\ker(S^*: \operatorname{Pic}^0(A)\to
 \operatorname{Pic}^0(A_d))\longrightarrow\mu_\infty.
\]
This is perfect: its character line bundles are exactly the summands
of $S_*\mathcal O_{A_d}$; see Milne, \emph{Abelian Varieties},
version 2.00, Chapter I, Theorem 9.1, printed p.~41.
Restricting to $K$ is therefore surjective
onto $K^\vee$ on the deck-group side. Since $\mathcal E'\ne0$, all
sixteen character summands occur and deck translations permute them
transitively.

Let $F$ be the trivial-character summand. The action of $T_S$ on labels
is multiplication by its scalar commutator character with $K$. This
defines a homomorphism $T_S\to K^\vee$: the underlying group
$T_S\times K$ is abelian and the coherence commutators are scalar.
It is onto by the deck pairing. Its kernel
\[
 H_S=\ker(T_S\to K^\vee)
\]
is an actual subgroup of translations of $A_d$, of order $4096$.
The coherent $T_S$-action restricts to $F$ for $H_S$. This retains the
equivariant lift; no nontrivial kernel action has been discarded.

In each degree, restriction to one character block and finite-\'etale
descent give exact identifications
\begin{equation}\label{hd-n-character-ext}
 \operatorname{Ext}^j(\mathcal E',\mathcal E')^{T_S\times K}
 \simeq\operatorname{Ext}^j(F,F)^{H_S}
 \simeq\operatorname{Ext}^j(\mathcal E,\mathcal E)^G.
\end{equation}
For the first identification, $K$-invariance removes all off-diagonal
blocks, while transitivity identifies the sixteen remaining blocks
from one stabilizer-invariant block. For the second, first descend
through $C_S$, then take the old $G$ invariants. Pullback on $HH$ is an
isomorphism and intertwines the traces, so HD-N03 proves injectivity
of trace on the left in degree two. On a class induced from the
$F$ block, additivity and translation invariance give
\[
 \sigma_{\mathcal E'}=16\sigma_F.
\]
Hence $\sigma_F$ is injective on
$\operatorname{Ext}^2(F,F)^{H_S}$, proving strong and therefore weak
$H_S$-semiregularity.

There are no negative self-Ext groups. Projection formula expresses
$\operatorname{Ext}^j(\mathcal E',\mathcal E')$ as the direct sum of
$\operatorname{Ext}^j(\mathcal E,\mathcal E\otimes P_\gamma)$ over
the character line bundles in $S_*\mathcal O_{A_d}$. By Mukai's formula
(3.1), conjugating tensor $P_\gamma$ through the equivalence in
HD-N03 gives a translation/tensor functor. The two arguments therefore
become coherent sheaves, so these groups vanish for $j<0$. Passing to
the direct summand $F$ preserves the vanishing.

All sixteen summands have identical rational Chern characters, since
they are translates. Thus
\begin{equation}\label{hd-n-character-sixteenth}
 \operatorname{ch}(F)=\frac1{16}S^*\operatorname{ch}(\mathcal E),
 \qquad \operatorname{rank}(F)=4.
\end{equation}
Use integral generators $A_j,B_j,C_j,D_j$ on the standard pairs, with
\[
 \theta_0=\sum_j(A_jB_j+C_jD_j),\quad
 R_j=A_jD_j-B_jC_j,\quad Q_j=A_jC_j+B_jD_j,
\]
and define
\begin{equation}\label{hd-n-standard-weil}
 t_d=-\bigl(R_1R_2R_3-R_1Q_2Q_3-Q_1R_2Q_3-Q_1Q_2R_3\bigr).
\end{equation}
With $i_0^*A=-B$ and $i_0^*B=A$, the matrix in
\eqref{hd-n-S} gives
\[
 S^*\Omega_{s,j}=8(R_j-sQ_j),\quad S^*t=512t_d,\quad
 S^*\theta=8\theta_0,\quad B_d=S^*B=2\sum_jQ_j.
\]
The thirty-two monomials of $t_d$ have coefficients $\pm1$, so $t_d$
is primitive integral. It is a Weil class because $S$ intertwines the
specified actions. Projection formula gives
\[
 \int t_d^2=\frac{4096(-2048)}{512^2}=-32,\qquad
 \theta_0t_d=0.
\]
Combining \eqref{hd-n-original-character} with
\eqref{hd-n-character-sixteenth} yields
\[
 e^{B_d}\operatorname{ch}(F)
 =32t_d+4-8\theta_0^2+\frac83\theta_0^4-\frac{16}{45}\theta_0^6.
\]
Each $Q_j$ is the first Chern class of the Poincar\'e bundle pulled
back by $(\mathrm{id},i_0)$ on the corresponding pair. Let $N$ be the
product of their squares, so $c_1(N)=B_d$, and set
$\mathcal E_d=F\otimes N$.
For the convention $\phi_N(h)=t_h^*N\otimes N^{-1}$, conjugating
$t_h^*$ by tensor $N$ is exactly
$t_h^*(-)\otimes(N\otimes t_h^*N^{-1})$. Thus the finite coherent
group for $\mathcal E_d$ is the actual graph
\[
 G_d=\{(h,-\phi_N(h)):h\in H_S\}\subset A_d\times\widehat A_d,
 \qquad |G_d|=4096.
\]
Its translation coordinate proves injectivity. The equivalence preserves
the semiregularity and negative-Ext statements. Its full character is
\eqref{hd-n-final-character}, with $B_0=0$ and $(\delta,T)=(1,32)$.

\paragraph{Deformation consequence and quantitative scope.}
Perry, \emph{The semiregularity theorem for equivariant noncommutative
varieties}, arXiv:2604.00511v2, Theorem 1.2, p.~3, applies to this
actual object on the connected standard balanced Weil sixfold family
after a finite level marking fixes the flat Weil class. Its hypotheses
are the coherent finite subgroup, weak semiregularity, negative
self-Ext vanishing, and the entire flat Hodge character proved above.
The supplied DEG\_L S-14 supplies the bounded-division consequence.
This construction does not yet extract a numerical effective
Chow-degree bound for the deformed perfect complex.

\paragraph{Precise source and certificate record.}
The categorical trace is identified with the Buchweitz--Flenner map
by Perry, Remark 2.5, p.~5; functoriality is Remark 2.3, p.~5;
Definition 2.6, pp.~5--6, and Lemma 2.7, p.~6, give the equivariant
criterion used here. The direct trace/evaluation identity is
Buchweitz--Flenner, Proposition 6.2.1, Corollary 6.3.2 and Theorem 6.4.2
(author preprint arXiv:math/0606730, pp.~29--31). Mukai's exact
equivalence and functor identities are Theorem 2.2, p.~156, formula
(3.1), p.~158, and formula (3.4), p.~159, in the published 1981 paper.
The exact exterior-algebra certificate checks all four one-pair
Fourier integrals, the full twelve-generator character, both Weil
norms, the two polarization volumes, the isogeny matrices, the
pullback factor $512$, and the character-summand factor $1/16$.

\subsection{HD-N07: every imaginary quadratic field on the same CM power}
\label{hd-n07}

The construction also gives an actual semiregular complex for every
split imaginary quadratic Weil sixfold family. The underlying CM
variety remains $E_i^6$; its new rational $K$-action is specified below.
Markman's supplied preprint already states algebraicity of all such
split sixfolds. Here the ideal-curve and Fourier--Mukai design is
credited to that source, while the required perfect-complex and trace
proof is checked directly, without assuming its algebraicity theorem.

\paragraph{The bielliptic disjointness calculation.}
For one bielliptic quotient $f:C\to E_i$ with involution $\sigma$,
\begin{equation}\label{hd-n-bielliptic-intersection}
 f^*\operatorname{Pic}^0(E_i)\cap(C_0-C_0)=\{0\}.
\end{equation}
Indeed, $p-q=f^*L$ implies
$p+\sigma q\sim q+\sigma p$. On a nonhyperelliptic curve, two
distinct effective divisors of degree two cannot be linearly equivalent:
their ratio would supply a degree-two map to $\mathbb P^1$.
The unordered divisors therefore agree. If $p\ne q$, both are fixed
points of $\sigma$.

The canonical $\sigma$-linearization of $\mathcal O_C(p-q)$ has
fibre sign $-1$ at $p$ and $q$, and $+1$ at the other two fixed points.
For example, in a local parameter $z$ with $\sigma z=-z$, both
$z^{-1}$, the generator of $\mathcal O(p)$, and $z$, the generator of
$\mathcal O(-q)$, change sign. Any two linearizations differ by a
global character of $\langle\sigma\rangle$, since automorphisms of a
line bundle are nonzero constants. Neither possible global sign makes
all four inertia actions trivial. A pullback $f^*L$ has trivial
inertia at all fixed points, so cannot be this line bundle. This proves
\eqref{hd-n-bielliptic-intersection}.

\paragraph{Quotients of the needed orders.}
Let $d$ be a positive squarefree integer, put $s^2=-d$, $K=\mathbb Q(s)$,
and $k=d+1$. Choose $\tau\in E_i$ of exact order $k$ and let
$\epsilon=f^*\tau$. The injectivity in HD-N01 makes its order $k$.
Equation \eqref{hd-n-bielliptic-intersection} makes all curves
$C_0+j\epsilon$, $0\le j<k$, pairwise disjoint. Thus
\[
 u_k:J(C)\to Y_k=J(C)/\langle\epsilon\rangle
\]
embeds $C_0$ as a smooth curve $\overline C_k$. If $u_k^*h=\Theta$ and
$\overline F_k=I_{\overline C_k}$, then
\begin{equation}\label{hd-n-general-ideal}
 \operatorname{ch}(\overline F_k)
 =1-\frac{k}{2}h^2+\frac{k}{3}h^3.
\end{equation}
This is the pullback calculation for $k$ disjoint Abel curves.

The maps in HD-N02 give, without change,
\[h^0(N_{\overline C_k/Y_k})=3,\qquad
\operatorname{ext}^1(\overline F_k,\overline F_k)
=\operatorname{ext}^2(\overline F_k,\overline F_k)=6.\]
The trace calculation still has rank six: the rank term is one and
the coefficient $-k/2$ of $h^2$ is nonzero. The degree-one trace also
remains injective. Thus both $\overline F_k$ and its external square
have the semiregularity used in HD-N03.

Put $r_k=u_kq:X\to Y_k$, $F_k=r_k^*\overline F_k$, and
\[
 H_k=\ker r_k=\ker q+\langle(\tau,0,0)\rangle,
 \qquad |H_k|=8k.
\]
The two summands intersect trivially because $f^*$ is injective.
Canonical descent makes $F_k\boxtimes F_k$ strongly
$H_k^2$-semiregular.

\paragraph{The actual group and full character.}
With exactly the $L,\mu,\Phi$ of HD-N03, define
\[
 \mathcal E_K=\Phi((F_k\otimes L)\boxtimes(F_k\otimes L)).
\]
It is perfect and has no negative self-Ext. For the convention $t_x^*$
and the inverse kernel $\mathcal P^{-1}[3]$, the transported group map
is exactly
\begin{equation}\label{hd-n-general-group}
 (x,y)\longmapsto
 \bigl((x-y,2(x+y)),(-2x,y)\bigr)
 \in A\times\widehat A.
\end{equation}
To check signs, tensor $L$ conjugates translation $x$ to translation
$x$ with tensor $P_{-2x}$. Pullback by $\mu$ sends the translation
coordinates to $(x-y,y)$ and the tensor coordinates to
$(-2x,-2(x+y))$. The inverse Fourier kernel exchanges translation
$y$ with tensor $P_y$ and tensor $P_b$ with translation $-b$.
This proves \eqref{hd-n-general-group}. Its last coordinate recovers
$y$ and its first then recovers $x$, so it is injective. Its finite
image $G_K$ has order $64(d+1)^2$. Its full coherence is transported
from descent; trace functoriality proves strong $G_K$-semiregularity.

Since $r_k^*h=2H$, equation \eqref{hd-n-general-ideal} gives
\[
 \operatorname{ch}(F_k\otimes L)
 =1+2H-2dH^2-\frac{4d}{3}H^3
 =\alpha e^{2sH}+\beta e^{-2sH},\quad
 \alpha=\frac{1+1/s}{2},\quad\beta=\frac{1-1/s}{2}.
\]
On $A=X\times\widehat X$ put
\[
 \theta_d=4dH_x+H_{\widehat x},\quad B=-P/2,\quad
 \Omega_{s,j}=4da_jb_j+2sP_j-u_jv_j,\quad
 t_K=-\tfrac12\left(\prod_j\Omega_{s,j}
                         +\prod_j\Omega_{-s,j}\right).
\]
After multiplication by $e^{-P_{\rm out}/2}$, the one-pair
Fourier integrals of HD-N04 equal
\[4s\exp(s\theta_{d,\rm out}/(4d))\]
for equal signs and
$\Omega_s$ for opposite signs. Taking three factors gives the
equal-sign coefficients $32d\pm16(d-1)s$; the opposite-sign coefficient
is $\alpha\beta=(d+1)/(4d)$. Hence the entire character is
\begin{equation}\label{hd-n-general-character}
 e^B\operatorname{ch}(\mathcal E_K)
 =\frac{d+1}{2d}t_K+P_d(\theta_d),
\end{equation}
where
\begin{align*}
 P_d(\theta)={}&64d-8(d-1)\theta-2\theta^2
       +\frac{d-1}{12d}\theta^3+\frac{\theta^4}{96d}\\
 &\hspace{12mm}-\frac{d-1}{3840d^2}\theta^5
       -\frac{\theta^6}{46080d^2}.
\end{align*}
The symbolic parameter certificate
\texttt{cm\_all\_k\_character.py} checks all four Fourier integrals
and this full identity over $\mathbb Q(d,s)/(s^2+d)$.

\paragraph{The precise Weil embedding and split form.}
Define the action of $s$ on each elliptic pair by
\[
 J_d=\begin{pmatrix}0&1/2\\-2d&0\end{pmatrix}.
\]
It squares to $-d$, is skew-adjoint for $\theta_d$, and has
holomorphic signature $(3,3)$. The order $\mathbb Z[2s]$ acts by
actual endomorphisms. The polarization has type
$(1,1,1,4d,4d,4d)$ and volume $64d^3$. The equality
$\Omega_s=-(u-2sa)(v-2sb)$ identifies its conjugate products with the
two Weil determinant lines. Thus $t_K$ is a nonzero Weil class.
It is primitive integral: odd powers of $s$ cancel, even powers are
integral, and the coefficient of $\prod_ju_jv_j$ is $+1$.
The exact products
\[
 \Omega_s\Omega_{-s}=-16d\,abuv,
 \qquad \theta_dt_K=0,\qquad \int t_K^2=-2048d^3
\]
give $(\delta,T)=(1,2048d^3)$ and
$\lambda=(d+1)/(2d)$.

For an explicit split hermitian basis, use the first elliptic
$a$- and $b$-cycles $e,f$ in each pair as a $K$-basis. With the
alternating polarization form $E_\theta$, define
\[
 h(v,w)=E_\theta(v,J_dw)+sE_\theta(v,w),
\]
conjugate-linear in its first argument. Its matrix is
$\left(\begin{smallmatrix}0&4ds\\-4ds&0\end{smallmatrix}\right)$.
The basis $e,f/(4ds)$ gives the hyperbolic matrix $U$; three pairs give
$U^3$. This is the exact rational hermitian isometry. The underlying
CM product is $E_i^6$ with the displayed balanced $K$-embedding.
For another fixed split polarized lattice, choose the rational
$K$-hermitian isometry $T$ to this source and an integer $M$ with
$MT\Lambda_{\rm target}\subset\Lambda_{\rm CM}$. On their common
finite-level period family the complex-linear map $MT$ induces an
actual isogeny of corresponding fibres, with pullback polarization
$M^2\theta_{\rm target}$.

The determinant monodromy is an algebraic integer unit of $K$ of norm
one, hence a root of unity. Killing this finite character makes $t_K$
flat. Perry's Theorem 1.2 applies to the full flat Hodge expression
\eqref{hd-n-general-character}, with the actual coherent group,
negative-Ext vanishing and semiregularity proved above. It supplies
the algebraicity conclusion in the given DEG\_L S-14 framework for
these split sixfold families. HD-N08--HD-N09 below obtain an
effective degree for every class-number-one field through the
explicit modified construction with two different bielliptic
quotient directions, reflection, Fourier transform and derived
dual. Its projective-dimension-one sheaf presentation and
regularity560 give effective evaluation-kernel third Chern cycles
and the bound $10^{258}d^{69}f^{33}\Delta^6$.
Nonsplit sixfold determinant data remain outside the proved
range of this construction.
